\section{Часть 2. Управление доступом к приложениям и системам базам данных}
\label{sec:part2}

PostgreSQL~\mdash  объектно-реляционная система управления базами данных с развитой моделью ролей и привилегий (\texttt{GRANT}/\texttt{REVOKE}) на уровне баз данных, схем, таблиц, последовательностей и функций. Роли могут наследовать права других ролей, что позволяет отделить политику доступа (\texttt{*\_role}) от учётных записей входа (\texttt{*\_user}).

Скрипты части~2 расположены в каталоге \texttt{lab1/p2} (таблица~\ref{tab:p2-scripts}).

\begin{table}[ht]
	\begin{tabularx}{\textwidth}{|l|X|}
		\hline
		\textbf{Скрипт} & \textbf{Назначение} \\
		\hline
		\texttt{install.sh} & Установка пакетов PostgreSQL и запуск службы \\
		\hline
		\texttt{setup\_acl.sh} & Создание БД \texttt{lab\_db}, схемы, ролей, таблицы и настройка \texttt{pg\_hba.conf} \\
		\hline
		\texttt{verify.sh} & Матрица операций SELECT / INSERT / UPDATE / DELETE / CREATE / DROP \\
		\hline
		\texttt{cleanup.sh} & Удаление объектов лабы; ключ \texttt{-{}-purge} удаляет пакеты СУБД \\
		\hline
	\end{tabularx}
	\caption{Скрипты части 2}
	\label{tab:p2-scripts}
\end{table}

В системе присутствуют три роли в соответствии с методичкой (таблица~\ref{tab:p2-roles}). Объекты стенда: база данных \texttt{lab\_db}, схема \texttt{lab\_schema}, таблица \texttt{lab\_schema.test\_table}.

\begin{table}[ht]
	\begin{tabularx}{\textwidth}{|l|l|l|X|}
		\hline
		\textbf{Учётка} & \textbf{Пароль} & \textbf{Роль} & \textbf{Права} \\
		\hline
		\texttt{admin\_user} & \texttt{admin\_pass} & \texttt{admin\_role} & полный доступ к схеме (DML и DDL) \\
		\hline
		\texttt{app\_user} & \texttt{user\_pass} & \texttt{user\_role} & SELECT, INSERT, UPDATE, DELETE \\
		\hline
		\texttt{guest\_user} & \texttt{guest\_pass} & \texttt{guest\_role} & только SELECT \\
		\hline
	\end{tabularx}
	\caption{Роли и учётные записи PostgreSQL}
	\label{tab:p2-roles}
\end{table}

Порядок запуска:

\begin{lstlisting}[caption={Запуск сценариев части 2},label={lst:p2-run}]
cd lab1/p2
chmod +x *.sh
sudo ./install.sh
sudo ./setup_acl.sh
./verify.sh
# sudo ./cleanup.sh
\end{lstlisting}

Прогон выполнен 10.09.2026. Проверка политики безопасности завершилась без расхождений с ожиданиями (\texttt{FAILS=0}). Сводка~\mdash  таблица~\ref{tab:p2-res}; полный лог~\mdash  \texttt{lab1/p2/results/verify\_acl.txt}.

\begin{table}[ht]
	\begin{tabularx}{\textwidth}{|l|c|c|c|c|c|c|}
		\hline
		\textbf{Учётка} & SELECT & INSERT & UPDATE & DELETE & CREATE & DROP \\
		\hline
		\texttt{admin\_user} & ДА & ДА & ДА & ДА & ДА & ДА \\
		\hline
		\texttt{app\_user} & ДА & ДА & ДА & ДА & НЕТ & НЕТ \\
		\hline
		\texttt{guest\_user} & ДА & НЕТ & НЕТ & НЕТ & НЕТ & --- \\
		\hline
	\end{tabularx}
	\caption{Результаты проверки ACL PostgreSQL, 10.09.2026}
	\label{tab:p2-res}
\end{table}

Таким образом, администратор имеет полный доступ, пользователь приложения~\mdash  только DML (CRUD), гость~\mdash  только чтение, что соответствует требованиям методички. Фрагмент назначения привилегий приведён в листинге~\ref{lst:p2-grant}.

\begin{lstlisting}[caption={Назначение привилегий ролям},label={lst:p2-grant}]
GRANT ALL PRIVILEGES ON SCHEMA lab_schema TO admin_role;
GRANT CREATE ON SCHEMA lab_schema TO admin_role;
GRANT SELECT, INSERT, UPDATE, DELETE
  ON ALL TABLES IN SCHEMA lab_schema TO user_role;
GRANT SELECT ON ALL TABLES IN SCHEMA lab_schema TO guest_role;
\end{lstlisting}
